\documentclass[]{article}

\usepackage{amsmath}
\usepackage{multirow}
\usepackage{natbib}
\usepackage{graphicx} 


\begin{document}

This study combines published observational data for two distinct populations of globular clusters (GCs) with a theoretical model of GC formation. Our methodology is designed to place both populations onto a common temporal framework to test the universality of the atomic cooling threshold as a primary regulator of their formation. All calculations are performed assuming a flat $\Lambda$CDM cosmology consistent with the Planck 2018 results: $H_0 = 67.4$ km s$^{-1}$ Mpc$^{-1}$, $\Omega_m = 0.315$, and $\Omega_\Lambda = 0.685$.

\subsection{Globular cluster samples}
We utilize data for two systems of high-redshift GCs. The first is the GEMS system, comprising 19 proto-GC candidates observed at a redshift of $z_{obs} = 9.625$. The second is the Sparkler system, which consists of 5 GCs observed at $z_{obs} = 1.4$. For each of the 24 clusters across both samples, we use the published stellar mass ($M_\star$), age, and metallicity ($[\mathrm{Z/H}]$) as the basis for our analysis. These properties were derived from detailed spectral energy distribution fitting in their respective discovery papers.

\subsection{Calculation of formation redshifts}
A central component of our analysis is the conversion of each cluster's measured age and observed redshift into a formation redshift, $z_{form}$. This places all clusters, regardless of their observation epoch, onto a common timeline of cosmic history. The formation redshift is defined as the redshift at which the age of the Universe was equal to the lookback time to the cluster's formation. We calculate $z_{form}$ by numerically solving the following equation for each cluster:
\begin{equation}
    t_L(z_{form}) = t_L(z_{obs}) + \tau
\end{equation}
where $\tau$ is the published age of the cluster and $t_L(z)$ is the lookback time to a given redshift $z$, calculated as:
\begin{equation}
    t_L(z) = \frac{1}{H_0} \int_0^z \frac{dz'}{(1+z')\sqrt{\Omega_m(1+z')^3 + \Omega_\Lambda}}
\end{equation}
The asymmetric uncertainties on the published ages were propagated through this calculation to determine the corresponding asymmetric uncertainties on $z_{form}$.

\subsection{Theoretical model of gc formation}
To provide a theoretical context for our empirical results, we compare the distribution of formation redshifts to the cosmic GC formation rate predicted by the atomic cooling threshold model. This model posits that the formation of the first GCs is restricted to dark matter halos with a virial temperature $T_{vir} \geq 10^4$ K, the threshold at which gas can cool efficiently via atomic hydrogen line emission. We computed the theoretical GC formation rate as a function of redshift by calculating the abundance of these halos using the Sheth-Tormen halo mass function, implemented with the same Planck 2018 cosmological parameters used throughout our analysis. The resulting curve represents the predicted rate at which GCs form within halos massive enough to support atomic cooling across cosmic time.

\subsection{Statistical analysis}
We characterized the properties of the GEMS and Sparkler populations by calculating descriptive statistics (mean, median, and standard deviation) for their stellar masses, ages, metallicities, and derived formation redshifts. To investigate the relationship between stellar mass and chemical enrichment, we analyzed the mass-metallicity relation (MZR) for each system individually and for the combined sample. We performed a robust linear regression on the $[\mathrm{Z/H}]$ versus $\log_{10}(M_\star/M_\odot)$ data using the Theil-Sen estimator. This non-parametric method is insensitive to outliers and provides a reliable measure of the underlying trend in the data. The statistical significance of the resulting MZR slope was assessed using its 95\% confidence interval.

\end{document}
                